\documentclass[10pt]{article}
% Draft formatted to Springer proceedings geometry; transfer to the official
% Complex Networks 2026 Overleaf template before submission (same structure).
\usepackage[a4paper,top=2.7cm,bottom=2.9cm,left=3.1cm,right=3.1cm]{geometry}
\usepackage{graphicx,booktabs,amsmath}
\usepackage[small,bf]{caption}
\usepackage{times}
\usepackage[numbers,sort&compress]{natbib}
\usepackage[hidelinks]{hyperref}
\setlength{\parskip}{1pt}
\captionsetup{font=footnotesize}

\title{\bfseries Steering a Network with Its Own Reflection:\\
True, False, and Noise Self-Model Broadcasts\\ in an Evolving Multi-Agent Network}
\author{Ivan Bhargava\\[2pt]\small Independent Researcher\\
\small \texttt{github.com/Ivan825/Project-ONE}}
\date{}

\begin{document}
\maketitle

\begin{abstract}
Many real networks respond to measurements of their own state: markets react to
indices, organizations to metrics, online platforms to trending signals. We
introduce a controlled laboratory for this reflexive loop. A population of
temporary agents harvests resources, cooperates, reproduces with mutation, and
dies, collectively forming an adaptive network $G(t)$. A global observer
periodically compresses the macrostate (fragmentation, centralization,
cooperation, inequality, turnover) into a self-model $S(t)$ that is broadcast
back to the agents either accurately, systematically distorted, or replaced by
matched-bandwidth noise; two further conditions withhold the broadcast, one
while still measuring. Across 970 deterministic, seed-replayable runs with
paired seeds and pre-registered outcomes, we find: (i) measurement without
broadcast is causally inert (outcome distributions identical to baseline,
Cliff's $\delta=0$); (ii) a systematically false self-model pulls the
macrostate toward itself far more strongly than noise of identical bandwidth
($\delta$ up to $0.85$, $p\!\approx\!10^{-8}$), with the effect largest, relative
to noise, at low feedback gain; (iii) the direction of the effect depends on
the type of distortion --- an inverted self-model exploits corrective
responses, a permanent-crisis self-model is self-defeating yet maximizes
cooperation, and a flattering self-model is behaviorally inert; and (iv)
heritable sensitivity to the broadcast is selected against in proportion to its
unreliability and volume: populations fed false self-models evolve to ignore
their observer within tens of generations. Accurate feedback roughly doubles
costly cooperation, while post-shock recovery time is insensitive to feedback
in all conditions --- a robust negative result. The framework offers a
reproducible causal testbed for reflexivity in adaptive networks, from
telemetry-driven infrastructure to measured social systems.
\end{abstract}

\section{Introduction}
Adaptive networks --- graphs whose topology coevolves with the dynamical states
of their nodes --- are a central object of network science
\cite{gross2008adaptive,holme2012temporal}. A distinctive subclass has received
far less controlled attention: networks that respond to \emph{descriptions of
themselves}. Financial markets move on published indices computed from their
own trades \cite{soros2013fallibility}; universities restructure around
rankings derived from their own behavior \cite{espeland2007rankings};
recommendation platforms amplify items \emph{because} an internal measurement
declared them trending. In each case a macroscopic measurement of the system is
compressed, published, and re-enters the system as an input to microscopic
decisions --- a loop Merton called the self-fulfilling prophecy
\cite{merton1948self} and whose measure-distorting variant is known as
Goodhart's law \cite{strathern1997improving}.

Empirical study of this loop is confounded: in real systems one can rarely
withhold the measurement, falsify it in controlled directions, or replay
history under identical randomness. We therefore built a minimal simulated
world in which the loop can be opened, closed, and corrupted at will, and
where every run is exactly reproducible from a configuration and seed. Agents
are deliberately simple --- transparent trait vectors rather than learned
policies --- so that every behavioral pathway from broadcast to response is
inspectable.

Our contributions are: (1) an open, deterministic experimental platform
coupling an evolving agent ecology, an adaptive interaction network, a global
observer, and a controllable broadcast channel; (2) a five-condition causal
design that separates the effect of \emph{being measured} from the effect of
\emph{being told}, and the effect of \emph{self-information} from that of
\emph{stimulation} (via a matched-bandwidth noise control); (3) evidence,
robust across feedback gains and two shock severities, that false self-models
steer the macrostate toward themselves while noise does not, and that the
steering direction is a function of the distortion type; and (4) to our
knowledge the first demonstration in an ecological agent-based setting that
\emph{attention to a collective self-model is itself an evolvable trait},
selected against in proportion to the self-model's unreliability.

\section{Related Work}
\textbf{Adaptive and temporal networks.} Coevolution of state and topology is
surveyed in \cite{gross2008adaptive,holme2012temporal}; our model belongs to
this family, with births, deaths and rewiring driven by agent decisions.
Robustness of networks to random and targeted node removal is classical
\cite{albert2000error}; we use targeted hub removal as a standardized
perturbation rather than an object of study.
\textbf{Global coupling in physics.} Globally coupled oscillators and maps
show that a mean field fed back to units reshapes collective dynamics
\cite{kuramoto1984chemical}. Our design differs in two ways central to the
research question: the ``mean field'' is a multidimensional, potentially
\emph{false} description, and units are born, die, and evolve.
\textbf{Reflexivity and reactive measurement.} Self-fulfilling prophecy
\cite{merton1948self}, reflexivity in markets \cite{soros2013fallibility}, and
the reactivity of rankings and metrics
\cite{espeland2007rankings,strathern1997improving} motivate our false-broadcast
conditions; we contribute a platform where these effects are dosable and
ablatable.
\textbf{Signaling and the evolution of trust.} In sender--receiver games,
receivers can evolve to ignore unreliable signals \cite{skyrms2010signals}.
Our distrust result is consonant with that theory but arises in an open
ecological population where the signal is \emph{self-referential} ---
a compressed description of the receiving population itself --- and where
attention to it is a continuous heritable trait under selection from an
energy economy, not a payoff matrix.
\textbf{Agent-based modeling.} Our methodology follows the generative ABM
tradition \cite{epstein1996growing,bonabeau2002agent}; the engine uses ideas
from Mesa \cite{kazil2020mesa} and NetworkX \cite{hagberg2008exploring}.
Multi-agent reinforcement learning studies global state as an
\emph{architectural} aid \cite{foerster2016learning}; we instead treat global
information as an \emph{experimental variable} with controlled veracity.
Populations of language-model agents are now deployed with shared dashboards
describing their own collective state \cite{park2023generative}; the present
rule-based setting provides causally clean baselines for that emerging regime.

\section{The Model}
\subsection{Agents and ecology}
Each agent $i$ carries a heritable trait vector
$\theta_i=(\mathrm{coop},\mathrm{expl},\mathrm{risk},\mathrm{soc},
\mathrm{share},\gamma_i)\in[0,1]^6$, where $\gamma_i$ is its
\emph{global sensitivity} --- how strongly broadcasts modulate its behavior.
Agents hold energy; each step they pay metabolism and choose one action ---
harvest, share (costly transfer to a neighbor), connect, prune, reproduce, or
idle --- sampled with probabilities proportional to trait- and state-dependent
propensities. The resource pool regenerates logistically, giving
density-dependent competition. Reproduction requires an energy threshold and
copies traits with Gaussian mutation ($\sigma=0.05$); energy exhaustion or age
causes death. Founding ages are staggered to avoid cohort artifacts.

\subsection{Observer and broadcast conditions}
Every $\Delta t=10$ steps the observer computes the macrostate $S(t)$:
population, fragmentation (1 minus largest-component share), degree- and
betweenness-concentration, Freeman centralization \cite{freeman1978centrality},
per-capita costly-helping rate, energy Gini, and turnover, plus mean trait
values (logged, never broadcast). Five conditions define what agents receive
(Table~\ref{tab:conditions}). Distorted broadcasts come in three modes:
\emph{invert} ($b_k = 1-s_k$ per normalized metric), \emph{crisis} (fixed
alarming values), and \emph{utopia} (fixed reassuring values). Broadcast
values shift action propensities through fixed corrective response rules ---
e.g., reported fragmentation above $0.3$ raises the propensity to form links;
a reported cooperation deficit raises the propensity to share --- each scaled
by $\gamma_i$ and a global gain $g$.

\begin{table}[t]
\centering\footnotesize
\caption{Experimental conditions. All start from matched parameter
distributions with paired random seeds.}
\label{tab:conditions}
\begin{tabular}{llll}
\toprule
Condition & Observer & Agents receive & Role\\
\midrule
A --- Local only & off & nothing & baseline\\
B --- Observed, blind & on & nothing & controls for measurement\\
C --- True feedback & on & accurate $S(t)$ & main treatment\\
F --- False feedback & on & distorted $S(t)$ & reflexivity probe\\
N --- Noise feedback & on & uniform random, same keys/ranges & stimulation control\\
\bottomrule
\end{tabular}
\end{table}

\subsection{Reproducibility}
The simulation is deterministic given (configuration, seed): one PRNG drives
all choices, agents act in sorted order, and a state hash certifies replay.
Identical hashes were verified across operating systems and machines. All
code, configurations, and analysis scripts are public; every figure regenerates
from two commands.

\section{Experimental Design}
Four outcomes were fixed before the first campaign: (1) post-shock recovery
time (steps until population regains 90\% of its pre-shock mean after
standardized hub removal); (2) mean post-shock fragmentation; (3) per-capita
costly-helping rate; (4) \emph{story pull} --- the mean signed movement of the
macrostate toward the broadcast, counted only where broadcast and reality
differ by more than $0.02$, so that the measure is well defined under false and
noise broadcasts. Campaigns: a \emph{flagship} run (conditions
A/B/C/F$_{\mathrm{invert}}$/N $\times$ 50 paired seeds, mild shock), a
\emph{harsh-shock} replication (removal of 40\% of hubs, $\times$30 seeds), and
a \emph{sensitivity sweep} (feedback gain $g\in\{0.2,0.8,1.6\}$ $\times$
\{A, C, F$_{\mathrm{invert}}$, F$_{\mathrm{crisis}}$, F$_{\mathrm{utopia}}$,
N\} $\times$ 20 seeds): 970 runs of 4{,}000 steps. Group differences use
Mann--Whitney $U$ with Cliff's $\delta$ effect sizes; all other logged
measures are exploratory.

\section{Results}
\subsection{Measurement alone is causally inert}
Under paired seeds, condition B is trajectory-identical to A: $\delta=0.000$
on every outcome in both the flagship and harsh-shock campaigns. The
instrument acquires causal power only when it broadcasts --- a validation that
the design isolates the channel of interest.

\subsection{A false self-model pulls the macrostate toward itself}
Story pull under the inverted false broadcast exceeds the noise control at
every gain tested (flagship: median $0.0067$ vs $0.0018$, $\delta=0.76$,
$p\!\approx\!6\times10^{-11}$; harsh shock: $0.0078$ vs $0.0021$,
$\delta=0.85$, $p\!\approx\!1.6\times10^{-8}$). Relative to noise the effect
is strongest at \emph{low} gain (ratio $\approx 7$ at $g=0.2$): a quiet,
systematic lie out-steers loud noise, while at high gain noise itself begins
to steer (Fig.~\ref{fig:sweep}, middle). Spillovers are substantial: the
inverted broadcast roughly halves post-shock fragmentation relative to
baseline ($\delta=-0.50$, $p\!\approx\!2\times10^{-5}$) because agents
persistently told the network is fragmenting over-connect it.

\begin{figure}[t]
\centering
\includegraphics[width=\textwidth]{sweep_summary.png}
\caption{Sensitivity sweep (40\% hub-removal shock; 20 seeds per cell;
medians). Left: evolved change in mean global sensitivity $\gamma$ ---
attention to the self-model is selected against in proportion to broadcast
unreliability and gain; the flattering (utopia) lie tracks the no-feedback
baseline exactly. Middle: story pull toward the broadcast. Right: cooperation
rate.}
\label{fig:sweep}
\end{figure}

\subsection{The type of lie sets the direction}
The three distortion modes produce qualitatively different regimes
(Fig.~\ref{fig:sweep}). The \emph{inverted} self-model, always reporting the
mirror of reality, generates the strongest pull: it keeps every corrective
response permanently active and thereby parasitizes the population's repair
machinery. The \emph{crisis} self-model is self-defeating for structure ---
corrections push fragmentation away from the alarmist report --- yet its
permanent cooperation-deficit signal mobilizes the highest costly-helping
rates observed ($0.36$--$0.45$). The \emph{utopia} self-model is behaviorally
inert: because responses fire only on reported deficits, a flattering
broadcast triggers nothing and is indistinguishable from disconnecting the
channel, at every gain. Under purely corrective response rules, flattery
equals silence.

\subsection{Populations fed lies evolve to stop listening}
Mean global sensitivity $\bar\gamma$ changes over $\sim$40 generations as an
emergent, unprogrammed function of broadcast quality
(Fig.~\ref{fig:trait}). At $g=0.8$ (harsh-shock campaign): no broadcast
$-0.03$ (neutral drift), truth $-0.10$, noise $-0.30$ ($\delta=-0.77$), false
$-0.42$ ($\delta=-0.94$, $p\!\approx\!4\times10^{-10}$). The sweep reveals a
dose--response: distrust of the false broadcast deepens from $-0.11$
($g=0.2$) to $-0.45$ ($g=1.6$) --- and at the highest gain even \emph{accurate}
feedback is selected against ($-0.24$): a broadcast that perturbs behavior
strongly becomes a fitness liability even when true. The utopia lie, being
consequence-free, breeds no distrust at all.

\begin{figure}[t]
\centering
\includegraphics[width=0.72\textwidth]{fig8_trait_evolution.png}
\caption{Evolution of attention to the self-model (harsh-shock campaign,
$g=0.8$, 30 seeds per condition). Change in population-mean global sensitivity
$\bar\gamma$ between the start and end of the run. The ordering
(false $<$ noise $<$ true $<$ none) emerged from selection alone.}
\label{fig:trait}
\end{figure}

\subsection{Accurate feedback as a distributed homeostat}
True feedback roughly doubles costly helping (flagship: $0.397$ vs $0.179$,
$\delta=0.96$, $p\!\approx\!10^{-16}$; replicated under harsh shock,
$\delta=0.95$), rising monotonically with gain ($0.27\!\to\!0.49$). The loop
functions as a population-level thermostat: reported deficits are repaired by
distributed individual action.

\subsection{A robust negative result}
Post-shock recovery time was insensitive to feedback in every condition and
campaign, even with 40\% of hubs removed (all $p>0.12$; median $\approx$
40 steps). In this ecology, resilience is demographic --- fast regrowth ---
rather than informational. We report this as a boundary of the phenomenon
rather than adjusting the protocol until the null disappears.

\section{Discussion and Limitations}
The first-order \emph{directions} of response are encoded in the agents'
corrective rules; that is by design, and we do not claim them as discoveries.
The findings live one level up: the magnitude and asymmetry of steering
(consistent lies $\gg$ matched noise, strongest at low gain), the spillovers
to non-targeted structure, the qualitative taxonomy of lies
(exploitative / mobilizing-but-self-defeating / inert), and above all the
evolutionary layer, in which the credibility of a collective self-model
becomes a selected property of the population. A conformist (non-corrective)
response variant is a natural next ablation and we expect it to reverse parts
of the taxonomy; the platform makes such tests one-line configuration changes.
Other limitations: a single ecology and parameter family (though effects are
monotone across an $8\times$ gain range), rule-based rather than learned
agents, and one network per run rather than interacting populations.

The results have a practical reading for engineered systems. Metric-driven
infrastructure --- autoscalers, observability pipelines, trending algorithms
--- constitutes exactly this loop; the false-broadcast conditions model
telemetry corruption, and the utopia result implies that \emph{flattering
telemetry is operationally equivalent to no telemetry} while raising no
alarms. For multi-agent AI, where fleets of agents increasingly receive
dashboards describing their own collective state, the evolution-of-distrust
result suggests that persistent misreporting does not merely misdirect a
population --- it degrades the control channel itself.

\section{Conclusion}
We introduced a reproducible causal laboratory for reflexive self-modeling in
adaptive networks and used it to show that (i) being measured changes nothing
while being told changes much; (ii) a consistent false self-description
steers a collective toward itself far more effectively than noise; (iii) the
kind of lie determines whether it exploits, mobilizes, or pacifies; and (iv)
selection acting on a heritable attention trait makes a population's trust in
its own self-model an evolving quantity --- one that unreliable observers
spend down. Future work extends the platform to conformist response rules,
continuity under total population turnover, and learned (including
language-model) agents.

\medskip
\noindent\textbf{Reproducibility.} Code, configurations, all 970 run
manifests, and one-command reproduction of every figure:
\texttt{github.com/Ivan825/Project-ONE}.

\begin{thebibliography}{25}\footnotesize
\bibitem{gross2008adaptive} Gross, T., Blasius, B.: Adaptive coevolutionary
networks: a review. J. R. Soc. Interface 5(20), 259--271 (2008)
\bibitem{holme2012temporal} Holme, P., Saram\"aki, J.: Temporal networks.
Phys. Rep. 519(3), 97--125 (2012)
\bibitem{soros2013fallibility} Soros, G.: Fallibility, reflexivity, and the
human uncertainty principle. J. Econ. Methodol. 20(4), 309--329 (2013)
\bibitem{espeland2007rankings} Espeland, W.N., Sauder, M.: Rankings and
reactivity: how public measures recreate social worlds. Am. J. Sociol.
113(1), 1--40 (2007)
\bibitem{merton1948self} Merton, R.K.: The self-fulfilling prophecy. Antioch
Rev. 8(2), 193--210 (1948)
\bibitem{strathern1997improving} Strathern, M.: `Improving ratings': audit in
the British university system. Eur. Rev. 5(3), 305--321 (1997)
\bibitem{albert2000error} Albert, R., Jeong, H., Barab\'asi, A.-L.: Error and
attack tolerance of complex networks. Nature 406, 378--382 (2000)
\bibitem{kuramoto1984chemical} Kuramoto, Y.: Chemical Oscillations, Waves,
and Turbulence. Springer (1984)
\bibitem{skyrms2010signals} Skyrms, B.: Signals: Evolution, Learning, and
Information. Oxford University Press (2010)
\bibitem{epstein1996growing} Epstein, J.M., Axtell, R.: Growing Artificial
Societies: Social Science from the Bottom Up. Brookings/MIT Press (1996)
\bibitem{bonabeau2002agent} Bonabeau, E.: Agent-based modeling: methods and
techniques for simulating human systems. PNAS 99(suppl.~3), 7280--7287 (2002)
\bibitem{kazil2020mesa} Kazil, J., Masad, D., Crooks, A.: Utilizing Python
for agent-based modeling: the Mesa framework. In: SBP-BRiMS. Springer (2020)
\bibitem{hagberg2008exploring} Hagberg, A., Schult, D., Swart, P.: Exploring
network structure, dynamics, and function using NetworkX. In: SciPy (2008)
\bibitem{foerster2016learning} Foerster, J., Assael, Y.M., de~Freitas, N.,
Whiteson, S.: Learning to communicate with deep multi-agent reinforcement
learning. In: NeurIPS (2016)
\bibitem{park2023generative} Park, J.S., O'Brien, J., Cai, C.J., Morris,
M.R., Liang, P., Bernstein, M.S.: Generative agents: interactive simulacra of
human behavior. In: UIST (2023)
\bibitem{freeman1978centrality} Freeman, L.C.: Centrality in social networks:
conceptual clarification. Soc. Netw. 1(3), 215--239 (1978)
\bibitem{gershenson2012complexity} Gershenson, C., Fern\'andez, N.:
Complexity and information: measuring emergence, self-organization, and
homeostasis at multiple scales. Complexity 18(2), 29--44 (2012)
\bibitem{hoel2013quantifying} Hoel, E.P., Albantakis, L., Tononi, G.:
Quantifying causal emergence shows that macro can beat micro. PNAS 110(49),
19790--19795 (2013)
\end{thebibliography}

\end{document}
